\chapter{Autoregressive Generative Models}

\section{Introduction}

Autoregressive models provide a simple and powerful approach to generative modeling by factorizing a joint distribution into a product of conditional distributions. They form the foundation of modern language models and connect naturally to sequence modeling, attention, and transformers.

This chapter develops autoregressive models from first principles, covering likelihood-based training, causal masking, and decoding strategies. We also discuss their role in modern large language models.

\section{Probability Factorization by the Chain Rule}

Any joint distribution over variables $x = (x_1,\dots,x_n)$ can be factorized as:
\[
p(x) = \prod_{i=1}^n p(x_i \mid x_{<i}),
\quad \text{where } x_{<i} = (x_1,\dots,x_{i-1}).
\]

\begin{itemize}
    \item this is an exact identity,
    \item no approximation is introduced,
    \item modeling reduces to learning conditional distributions.
\end{itemize}

\section{Autoregressive Models for Different Data Types}

\subsection{Sequences (Text)}

For text:
\[
p(x_1,\dots,x_T) = \prod_{t=1}^T p(x_t \mid x_{<t}).
\]

Each token is predicted from previous tokens.

\subsection{Images}

Images can be flattened into sequences (e.g., row-wise):
\[
p(x) = \prod_{i=1}^N p(x_i \mid x_{<i}).
\]

Here $x_i$ corresponds to pixel intensities.

\subsection{General Data}

Any structured data can be ordered and modeled autoregressively, though the choice of ordering affects modeling efficiency.

\section{Log-Likelihood Objective}

Given data $\{x^{(j)}\}$, the log-likelihood is:
\[
\sum_j \log p(x^{(j)}) = \sum_j \sum_{i=1}^n \log p(x_i^{(j)} \mid x_{<i}^{(j)}).
\]

\begin{itemize}
    \item training decomposes into local prediction problems,
    \item likelihood is tractable and exact,
    \item gradients are straightforward to compute.
\end{itemize}

\section{Teacher Forcing and Next-Token Prediction}

During training, models are typically trained with \emph{teacher forcing}:
\begin{itemize}
    \item the true previous tokens $x_{<t}$ are provided,
    \item the model predicts $x_t$.
\end{itemize}

This leads to the objective:
\[
\sum_t \log p(x_t \mid x_{<t}),
\]
which corresponds to next-token prediction.

\section{Causal Masking and Training}

In neural implementations (e.g., transformers), causal structure is enforced by masking.

\begin{definition}[Causal Masking]
During training, ensure that the prediction at position $i$ depends only on $x_{<i}$ by masking future positions:
\[
\text{attention weight } \alpha_{ij} = 0 \quad \text{for } j > i.
\]
\end{definition}

This ensures that the model respects the autoregressive factorization.

\section{Sampling and Decoding}

To generate samples, we proceed sequentially:
\[
x_1 \sim p(x_1),\quad
x_2 \sim p(x_2 \mid x_1),\quad
\dots
\]

\subsection{Decoding Strategies}

\begin{itemize}
    \item \textbf{Greedy decoding:} choose most probable token,
    \item \textbf{Sampling:} sample from the distribution,
    \item \textbf{Temperature scaling:} adjust randomness,
    \item \textbf{Top-$k$ / nucleus sampling:} restrict candidate tokens.
\end{itemize}

\subsection{Tradeoffs}

\begin{itemize}
    \item deterministic decoding yields coherent but less diverse outputs,
    \item stochastic decoding increases diversity but may reduce quality.
\end{itemize}

\section{Exposure Bias}

During training:
\begin{itemize}
    \item model sees true history,
\end{itemize}

During generation:
\begin{itemize}
    \item model conditions on its own predictions.
\end{itemize}

This mismatch is called \emph{exposure bias}.

\subsection{Implication}

Errors can accumulate during generation, as mistakes affect future predictions.

\section{Worked Examples}

\subsection{Character-Level Language Modeling}

For a text sequence:
\[
p(x) = \prod_{t=1}^T p(x_t \mid x_{<t}),
\]
where $x_t$ are characters.

The model predicts each character given previous ones.

\subsection{Pixel Ordering in Images}

Flatten an image into a sequence:
\[
(x_{1,1}, x_{1,2}, \dots, x_{H,W}).
\]

Then:
\[
p(x) = \prod_i p(x_i \mid x_{<i}).
\]

This models spatial dependencies sequentially.

\subsection{Simple Decoding Example}

Given probabilities:
\[
p(x_t \mid x_{<t}) = (0.7, 0.2, 0.1),
\]
\begin{itemize}
    \item greedy decoding selects the first token,
    \item sampling may choose any token according to probabilities.
\end{itemize}

\section{Relationship to Modern Language Models}

Modern large language models are autoregressive models trained on massive text corpora.

\begin{itemize}
    \item transformers implement $p(x_t \mid x_{<t})$,
    \item causal masking enforces autoregression,
    \item training is next-token prediction.
\end{itemize}

Thus autoregressive modeling is the foundation of modern NLP systems.

\section{Comparison with Latent-Variable Models}

\begin{itemize}
    \item autoregressive models:
    \begin{itemize}
        \item exact likelihood,
        \item sequential generation,
        \item no latent variables required.
    \end{itemize}
    \item latent-variable models:
    \begin{itemize}
        \item introduce hidden structure,
        \item may require approximate inference,
        \item allow global representations.
    \end{itemize}
\end{itemize}

Each approach has strengths and limitations.

\section{A Unifying Perspective}

Autoregressive models reduce generative modeling to conditional prediction. By factorizing distributions via the chain rule, they transform a complex global problem into a sequence of local prediction problems.

They connect naturally to:
\begin{itemize}
    \item sequence modeling,
    \item attention mechanisms,
    \item modern language models.
\end{itemize}

\section{Summary}

In this chapter we developed autoregressive generative models. We derived the chain-rule factorization, introduced likelihood-based training, and explained teacher forcing and causal masking. We discussed sampling and decoding strategies, as well as exposure bias.

Autoregressive models provide a simple yet powerful framework that underlies modern large-scale generative systems.

\section*{Exercises}

\begin{enumerate}
    \item Write the autoregressive factorization for:
    \begin{enumerate}
        \item a sequence of words,
        \item a grayscale image,
        \item a time series.
    \end{enumerate}

    \item Show that maximizing log-likelihood in an autoregressive model reduces to a sum of conditional log-probabilities.

    \item Explain the difference between teacher forcing and free-running generation.

    \item Compare greedy decoding and sampling. When might each be preferred?

    \item Explain exposure bias and propose one possible mitigation strategy.

    \item Compare autoregressive models and latent-variable models in terms of:
    \begin{enumerate}
        \item likelihood computation,
        \item sampling,
        \item representation learning.
    \end{enumerate}

    \item For a simple sequence model, compute $p(x)$ given conditional probabilities.

    \item Explain how causal masking enforces autoregressive structure in transformers.
\end{enumerate}